\subsection{The unramified double-cover repair and a genus-three obstruction}
\label{sec:genus-three-pointing-counterexample}

\paragraph{Finding A30 (the genus-two character truncation is essential).}
Over $k=\mathbf F_4$, let $D$ be the smooth projective model of
\[
                   z^2+z=\frac{1}{(t^2+t)^3}.
\tag{DC.1}
\]
It has genus three and $2$-rank one.  Let $P$ be the point
above $t=\infty$ with $z=0$, and let $W$ be the unique point
above $t=0$.  Both are rational; $W$ is Weierstrass and $P$
is not.  There is a surjection
\[
            b:J_D(k)\longrightarrow\mathbf Z/3,
                 \qquad b([W-P])\ne0,
\tag{DC.2}
\]
whose associated constant, geometrically connected, $P$-split
cyclic \'etale cover $D'\to D$ has \emph{even} relative class
number.  Since $\operatorname{ord}_3(2)=2$, the genus-two
conclusion of A28 that every such detecting quotient gives
an odd relative class number does not extend to genus three.
This does not rule out witnesses from other quotients or
higher powers for the same pointed genus-three curve.
The canonical class is also
detected: $K_D\sim4W$, so $b([K_D-4P])\ne0$.

More precisely, if $H_Y(T)$ denotes the numerator of the zeta
function of a curve $Y/k$, and
$B=(\ker(J_{D'}\to J_D))^0$, then
\[
 \overline{H_D}(T)=1-T,\qquad
 \overline{H_{D'}}(T)=(1-T)^2(1-T^3),\qquad
 \overline{H_B}(T)=(1-T)(1-T^3)
                         \quad\text{in }\mathbf F_2[T].
\tag{DC.3}
\]
The cover has degree three and $g(D')=7$.

\begin{proof}
\emph{1. The curve and its elliptic quotient.}
The rational function in (DC.1) has poles of order three
at $t=0,1$, and no others.  The Artin--Schreier genus and
$2$-rank formulas give $g(D)=3$ and $f_D=1$
\cite[Lemma~2.6, p.~712]{PriesZhu2012}.
Its hyperelliptic involution is $z\mapsto z+1$.
The two points above infinity have $z=0,1$ and are
unramified; the points above $0,1$ are its two ramification
points.  This proves the assertions about $P$ and $W$.
Since the hyperelliptic degree-two divisor class is $2W$,
the canonical class in genus three is $4W$.

Put
\[
 E/k:\ y^2+y=x^3,\qquad O=\infty,\qquad Q=(0,0).
\]
There is a separable degree-two map
\[
 f:D\longrightarrow E,\qquad
                 x=(t^2+t)^{-1},\quad y=z.
\tag{DC.4}
\]
Indeed its function-field extension is
$t^2+t=1/x$.  The two simple poles of $1/x$ on $E$ are
$Q=(0,0)$ and $-Q=(0,1)$, so these and only these are its
branch points.  Furthermore $f(P)=Q$ and $f(W)=O$.

There are nine points on $E(k)$: each of the four values
of $x$ has two values of $y$, and there is $O$.
Every affine point $(a,c)$ has order dividing three.
Its tangent line is $y=c+a^2(x-a)$; substituting
$x=a+v$ and using $a^4=a$ leaves $v^3=0$.
Thus the tangent meets the cubic three times at that point.
Consequently
\[
              E(k)\simeq(\mathbf Z/3)^2,
              \qquad \operatorname{ord}(Q)=3.
\tag{DC.5}
\]
The point count also gives
$H_E(T)=1+4T+4T^2$, in particular $\overline{H_E}=1$.

\emph{2. A translated Lang isogeny with the required fibres.}
Choose $a:E(k)\twoheadrightarrow\mathbf Z/3$ with $a(Q)=1$,
and write $H=\ker a$.
Let $F$ be $4$-power Frobenius on $E$.  Factoring the Lang
isogeny through $H$ gives
\[
 E\longrightarrow Y:=E/H\xrightarrow{\psi}E,
             \qquad \psi\pi_H=F-1.
\tag{DC.6}
\]
The map $\psi$ is an \'etale isogeny of degree three with
constant kernel $E(k)/H$.  The usual Lang-fibre computation
of A12 shows that a fibre $\psi^{-1}(R)$ is rational exactly
when $R\in H$; if rational, all three points are rational.
Use the translated cover
\[
                  \eta:Y\longrightarrow E,
                           V\longmapsto\psi(V)+Q.
\tag{DC.7}
\]
Its fibre above $Q$ has three rational points.  Its fibre
above $-Q$ has Frobenius translation
$a(-Q-Q)=a(Q)=1$, because $3Q=0$.
It is therefore one closed point of degree three.
The fibre above $O$ likewise has nontrivial translation
$a(-Q)=-1$.

Define $D'=D\times_{E,\eta}Y$.  This is a degree-three
constant \'etale cover of $D$, and its fibre above $P$ is
split.  It is geometrically connected: over an algebraic
closure the function-field extensions of $E$ defining $D$
and $Y$ have coprime degrees two and three, and the latter
is Galois.  They are linearly disjoint.
Unramified Riemann--Hurwitz gives $g(D')=3(3-1)+1=7$.

The rational quotient defining this pointed cover is
$b=a\circ f_*:J_D(k)\to\mathbf Z/3$.
Indeed $f_*[R-P]=[f(R)-Q]$, which is exactly the Lang
translation used in (DC.7).  In particular
\[
                 b([W-P])=a(O-Q)=-1\ne0.
\tag{DC.8}
\]
This also proves that $b$ is surjective.

\emph{3. The exact zeta calculation modulo two.}
For a separable degree-two cover $X\to Z$ of curves in
characteristic two, let $\mathcal R$ be its set of closed
branch points on $Z$.  Comparing Euler factors gives
\[
 Z_X(T)\equiv Z_Z(T^2)
                   \prod_{R\in\mathcal R}(1-T^{\deg R})
                                \pmod2.
\tag{DC.9}
\]
At an unramified closed point of degree $d$, either there
are two degree-$d$ points above it or one degree-$2d$ point.
Both local factors reduce to $(1-T^{2d})^{-1}$.
At a ramified point the local factor is $(1-T^d)^{-1}$;
multiplication by $1-T^d$ changes the corresponding
$Z_Z(T^2)$ factor to this one.  This proves (DC.9) as an
identity of formal power series.

Since the constant field has even cardinality, passing to
zeta numerators in (DC.9) gives the polynomial identity
\[
 \overline{H_X}(T)=\overline{H_Z}(T^2)
             \frac{\prod_{R\in\mathcal R}(1-T^{\deg R})}{1-T}.
\tag{DC.10}
\]
For $D\to E$, the two branch points are rational and
$\overline{H_E}=1$, giving $\overline{H_D}=1-T$.
The elliptic curve $Y$ is isogenous to $E$, so
$H_Y=H_E$ and $\overline{H_Y}=1$.
For $D'\to Y$, the branch divisor is the \'etale pullback
of $Q+(-Q)$.  Its closed-point degrees are $1,1,1,3$ by
Step~2.  Formula (DC.10) therefore gives
\[
       \overline{H_{D'}}(T)
              =\frac{(1-T)^3(1-T^3)}{1-T}
              =(1-T)^2(1-T^3).
\]
Pullback and norm for the degree-three cover show that
$J_D\times B$ is isogenous to $J_{D'}$.  Dividing their
Frobenius polynomials proves the last identity in (DC.3).
Finally
$\#B(k)=H_B(1)\equiv0\pmod2$.
The point-count and isogeny identities are the usual
Frobenius identities, as in
\cite[Theorem~19.1(b) and its Frobenius proof, pp.~40--41]{MilneAV}.

\emph{4. The individual reduced character polynomials.}
The failure is visible before taking their product.
Let $\chi:\mathbf Z/3\to\overline{\mathbf F}_2^\times$
send $1$ to $\zeta$, a primitive cube root.
The nontrivial unramified character sheaf on $E$ defined
by (DC.7) has $L$-polynomial $1$.
For example, Riemann--Roch on the elliptic curve makes
the number of effective divisors in a rational divisor
class of each positive degree independent of the class,
so every positive-degree nontrivial character coefficient
vanishes.

The character Euler-factor argument of Step~3 applies
with a character value $\xi$ at a closed point of $E$:
both unramified splitting types give
$(1-\xi^2T^{2d})^{-1}$ modulo two, while a branch point
contributes the correction $1-\xi T^d$.
The character values at $Q,-Q$ are $1,\zeta$ respectively.
Thus the degree-four unramified character polynomial on
the genus-three curve $D$ has reduction
\[
       \overline{L_{\chi,D,P}}(T)=(1-T)(1-\zeta T),
       \qquad
       \overline{L_{\chi^{-1},D,P}}(T)
                                      =(1-T)(1-\zeta^2T).
\tag{DC.11}
\]
Both values at $T=1$ vanish.  The genus-two reduction is
linear, whereas (DC.11) has a genuine quadratic term.
This is precisely why its inversion symmetry does not
give the genus-two contradiction.
\end{proof}

\paragraph{What happens to the proposed unramified double-cover repair.}
The underlying geometric observation is correct.
For a rank-one genus-two Artin--Schreier curve written,
after choosing the two rational branch points as $0,\infty$,
as
\[
 C:\ y^2+y=f_1(x)+f_3(x),
\]
where $f_1$ has just the simple pole at $0$ and $f_3$
just the triple pole at $\infty$, form the connected
biquadratic fibre product
\[
 z_1^2+z_1=f_1(x),\qquad z_2^2+z_2=f_3(x),
                         \qquad y=z_1+z_2.
\tag{DC.12}
\]
Over an algebraic closure the two Artin--Schreier classes
are independent, since their ramification supports differ.
The inertia groups over $0,\infty$ are respectively the
$z_1$- and $z_2$-translation groups.
The diagonal translation is in neither inertia group,
so its quotient map $\widetilde C\to C$ is \'etale of
degree two.  It is the unique nontrivial geometric etale
double cover: $f_C=1$ means that
$H^1_{\mathrm{et}}(C_{\bar k},\mathbf F_2)$ has dimension one.
Its source has genus three by Riemann--Hurwitz.
The quotient by $z_2$-translation is
the genus-zero $z_1$-curve.  Hence $z_2$-translation is
the hyperelliptic involution of $\widetilde C$.

At the simple pole its two lifts are interchanged by
this involution, so neither is Weierstrass.  At the
triple pole both lifts are fixed and are Weierstrass.
Equivalently, the hyperelliptic presentation on the
genus-zero $z_1$-curve has two poles of order three,
so the same formula gives $2$-rank one.
Finite constant extensions, or the appropriate constant
twist, can make a prescribed fibre rational, but do not
change these geometric conclusions.

This construction nevertheless does not repair the A1
argument.  First, its nontrivial geometric degree-two
quotient does not factor through the geometric
prime-to-two fundamental group, whose finite quotients
have odd order.  A collision in the quotient used by A1
therefore does not supply lifts to this cover.
Second, if $P_s,P_t$ are the simple- and triple-pole
points on $C$, their difference has exact order two:
$\operatorname{div}(x)=2P_s-2P_t$, and a nonzero
degree-one principal divisor would make $C$ rational.
For lifts $\widetilde P_s,\widetilde P_t$ one has
\[
 \pi_*[\widetilde P_t-\widetilde P_s]=[P_t-P_s].
\tag{DC.13}
\]
Every odd cyclic character pulled back from $C$ thus
kills this difference.  A new odd character detecting
it upstairs cannot supply the A28 detector by descent
from such a character on $C$.

The explicit curve (DC.1) is exactly this construction
for $C:y^2+y=x^3+x^{-1}$, with $z_1=t$ and $z_2=z$.
Its example (DC.11) gives an additional obstruction:
even as an argument about the upstairs curve alone,
the proposed unrestricted higher-genus symmetry rule
is false.
